\documentclass[11pt,a4paper]{article}

% =========================================================
% Computational Neuroscience Study Guide
% Author: Ali Taghizadeh
% =========================================================

% ---------- Page Layout ----------
\usepackage[margin=2.3cm]{geometry}
\usepackage[T1]{fontenc}
\usepackage{lmodern}
\usepackage{microtype}

% ---------- Mathematics ----------
\usepackage{amsmath, amssymb, mathtools}

% ---------- Tables and Lists ----------
\usepackage{booktabs}
\usepackage{array}
\usepackage{longtable}
\usepackage{enumitem}

% ---------- Colors and Boxes ----------
\usepackage{xcolor}
\usepackage{tcolorbox}
\tcbuselibrary{skins,breakable}

% ---------- Formatting ----------
\usepackage{titlesec}
\usepackage{fancyhdr}

% ---------- Hyperlinks ----------
\usepackage[hidelinks]{hyperref}

% =========================================================
% Colors
% =========================================================
\definecolor{navy}{RGB}{43,108,176}
\definecolor{burnt}{RGB}{221,107,32}
\definecolor{leafgreen}{RGB}{47,133,90}
\definecolor{lightgray}{RGB}{245,246,248}
\definecolor{softblue}{RGB}{235,244,255}
\definecolor{darkgray}{RGB}{74,85,104}
\definecolor{midgray}{RGB}{113,128,150}

% =========================================================
% Document Metadata
% =========================================================
\newcommand{\GuideTitle}{Computational Neuroscience Study Guide}
\newcommand{\GuideSubtitle}{From Passive Membranes to Balanced Spiking Networks}
\newcommand{\GuideSource}{Personal notes based on the \texttt{Neuro-AI-Foundations} repository}
\newcommand{\GuideAuthor}{Ali Taghizadeh}

\hypersetup{
	pdftitle={Computational Neuroscience Study Guide},
	pdfauthor={Ali Taghizadeh},
	pdfsubject={Computational Neuroscience and Spiking Neural Networks},
	pdfkeywords={computational neuroscience, LIF, EIF, AdEx, balanced networks, spiking neural networks}
}

% =========================================================
% Section Styling
% =========================================================
\titleformat{\section}
{\normalfont\Large\bfseries\color{navy}}
{\thesection}
{0.6em}
{}

\titleformat{\subsection}
{\normalfont\large\bfseries\color{darkgray}}
{\thesubsection}
{0.6em}
{}

\titlespacing*{\section}{0pt}{1.4em}{0.7em}
\titlespacing*{\subsection}{0pt}{1.1em}{0.45em}

% =========================================================
% Header / Footer
% =========================================================
\pagestyle{fancy}
\fancyhf{}
\renewcommand{\headrulewidth}{0.4pt}
\renewcommand{\footrulewidth}{0pt}

\fancyhead[L]{\small\textit{Neuro-AI-Foundations --- Study Guide}}
\fancyhead[R]{\small\GuideAuthor\quad|\quad\thepage}

% =========================================================
% List Styling
% =========================================================
\setlist[itemize]{topsep=3pt,itemsep=2pt,parsep=0pt,leftmargin=1.4em}
\setlist[enumerate]{topsep=3pt,itemsep=2pt,parsep=0pt,leftmargin=1.6em}

% =========================================================
% Custom Boxes
% =========================================================
\newtcolorbox{eqbox}[1][]{
	enhanced,
	colback=lightgray,
	colframe=navy,
	boxrule=0.6pt,
	arc=2pt,
	left=6pt,
	right=6pt,
	top=4pt,
	bottom=4pt,
	title={\textbf{Key Equation: #1}},
	fonttitle=\small\bfseries,
	coltitle=white,
	colbacktitle=navy,
	breakable
}

\newtcolorbox{keybox}[1][Key Takeaway]{
	enhanced,
	colback=green!5!white,
	colframe=leafgreen,
	boxrule=0.6pt,
	arc=2pt,
	left=6pt,
	right=6pt,
	top=4pt,
	bottom=4pt,
	title={\textbf{#1}},
	fonttitle=\small\bfseries,
	coltitle=white,
	colbacktitle=leafgreen,
	breakable
}

\newtcolorbox{defbox}[1][Definition]{
	enhanced,
	colback=orange!5!white,
	colframe=burnt,
	boxrule=0.6pt,
	arc=2pt,
	left=6pt,
	right=6pt,
	top=4pt,
	bottom=4pt,
	title={\textbf{#1}},
	fonttitle=\small\bfseries,
	coltitle=white,
	colbacktitle=burnt,
	breakable
}

\newtcolorbox{overviewbox}{
	enhanced,
	colback=softblue,
	colframe=navy,
	boxrule=0.7pt,
	arc=3pt,
	left=8pt,
	right=8pt,
	top=7pt,
	bottom=7pt,
	width=0.94\textwidth
}

% =========================================================
% Math Shortcuts
% =========================================================
\newcommand{\dd}{\mathrm{d}}
\newcommand{\e}{\mathrm{e}}

% =========================================================
% Title
% =========================================================
\title{}
\author{}
\date{}

\begin{document}
	
	% =========================================================
	% Cover Header
	% =========================================================
	\begin{center}
		\vspace*{-1.2cm}
		
		{\Huge\bfseries\color{navy}\GuideTitle\par}
		
		\vspace{0.25cm}
		
		{\Large\color{darkgray}\GuideSubtitle\par}
		
		\vspace{0.35cm}
		
		{\normalsize\itshape\color{midgray}\GuideSource\par}
		
		\vspace{0.45cm}
		
		{\large\bfseries Author: \GuideAuthor\par}
		

		
		\vspace{0.45cm}
		
		{\color{navy}\rule{0.82\textwidth}{0.8pt}}
		
		\vspace{0.55cm}
		
		\begin{overviewbox}
			\centering
			\small
			\textbf{Conceptual progression covered in this guide}\\[6pt]
			
			Passive membrane
			$\;\longrightarrow\;$
			Leaky Integrate-and-Fire, LIF
			$\;\longrightarrow\;$
			Exponential Integrate-and-Fire, EIF
			$\;\longrightarrow\;$
			Adaptive Exponential I\&F, AdEx
			$\;\longrightarrow\;$
			Balanced spiking networks
			
			\vspace{5pt}
			
			{\footnotesize
				\color{darkgray}
				no spikes
				\hspace{0.9cm}
				hard threshold
				\hspace{0.8cm}
				soft nonlinear spike onset
				\hspace{0.8cm}
				adaptation dynamics
				\hspace{0.8cm}
				large-scale network regime
			}
		\end{overviewbox}
		
		\vspace{0.45cm}
		
		\begin{tcolorbox}[
			enhanced,
			colback=lightgray,
			colframe=darkgray,
			width=0.94\textwidth,
			arc=3pt,
			boxrule=0.5pt,
			left=8pt,
			right=8pt,
			top=6pt,
			bottom=6pt
			]
			\small
			\textbf{Purpose of these notes.}
			This document provides a structured path through core computational neuroscience models,
			starting from the passive membrane equation and progressively building toward adaptive
			single-neuron models and balanced recurrent spiking networks. The emphasis is on
			intuition, mathematical form, and the relationship between biological mechanisms and
			AI-oriented network computation.
		\end{tcolorbox}
		
	\end{center}
	
	\vspace{0.2cm}
	
		\newpage
	\tableofcontents
	
	\newpage

% =====================================================================
\section{Unit Convention (read this first)}
% =====================================================================
Every model in this project uses one consistent unit system. Get comfortable with it before anything else, because it's what makes every equation below dimensionally consistent \emph{without} conversion factors.

\begin{center}
\begin{tabular}{@{}lll@{}}
\toprule
\textbf{Quantity} & \textbf{Symbol} & \textbf{Unit} \\
\midrule
Time & $t$ & ms \\
Voltage & $v$ & mV \\
Current & $I$ & pA \\
Capacitance & $C_m$ & pF \\
Conductance & $g_L$ & nS \\
Resistance & $R = 1/g_L$ & G$\Omega$ \\
\bottomrule
\end{tabular}
\end{center}

\begin{keybox}[Why this works]
$1\,\text{nS} \times 1\,\text{mV} = 1\,\text{pA}$, and $1\,\text{pF} \times 1\,\text{mV/ms} = 1\,\text{pA}$. So an equation like $C_m \dot v = -g_L(v-v_{rest}) + I$ needs no unit-conversion constants when everything is expressed in these units. This is the same convention used by Brian2 and NEST for these model families.
\end{keybox}

% =====================================================================
\section{The Passive (RC) Membrane}
% =====================================================================
The baseline model. A neuron's membrane behaves, at minimum, like a resistor and capacitor in parallel: charge leaks out through the resistor while the capacitor smooths sudden changes.

\begin{eqbox}[Passive membrane]
$$\tau_m \frac{\dd v}{\dd t} = -(v - v_{rest}) + R\,I(t) \qquad\Longleftrightarrow\qquad C_m \frac{\dd v}{\dd t} = -g_L(v - v_{rest}) + I(t)$$
where $\tau_m = C_m/g_L$ is the \textbf{membrane time constant}.
\end{eqbox}

\subsection{Step-current response (closed form)}
For a constant current $I$ switched on at $t=0$ with $v(0)=v_{rest}$:
$$v(t) = v_{rest} + \frac{I}{g_L}\left(1 - e^{-t/\tau_m}\right)$$
As $t \to \infty$, $v \to v_{rest} + I/g_L = v_{rest} + IR$ — the steady-state deflection is \emph{linear} in $I$. This is the definition of "passive": no voltage-dependent conductances, so the system is a linear low-pass filter with cutoff set by $\tau_m$.

\begin{keybox}
There is no threshold and no reset in this model. No matter how large $I$ is, $v(t)$ only asymptotes toward $v_{rest}+IR$ — it never spikes, because nothing in the equation \emph{can} produce a spike. This is the whole pedagogical point of including it: a threshold is not a detail, it's what makes ``firing rate'' a meaningful concept at all.
\end{keybox}

% =====================================================================
\section{Leaky Integrate-and-Fire (LIF)}
% =====================================================================
Keeps the passive equation exactly as-is, and bolts on the simplest possible spiking rule.

\begin{eqbox}[LIF spiking rule]
$$v(t) \geq v_{th} \;\Longrightarrow\; \text{spike recorded}, \quad v \leftarrow v_{reset}, \quad \text{then } v \text{ held for } t_{ref}\text{ ms}$$
\end{eqbox}

\subsection{Rheobase}
The \textbf{rheobase} is the minimum constant current needed for the neuron to ever reach threshold (i.e., for the steady state to reach $v_{th}$):
$$I_{rheo} = g_L\,(v_{th} - v_{rest})$$
Below $I_{rheo}$, the neuron behaves identically to the passive membrane and never spikes. Above it, it fires periodically forever.

\subsection{Analytical F-I curve}
For $I > I_{rheo}$ (and taking $v_{reset}=v_{rest}$ for simplicity), the time to charge from reset back up to threshold is $\tau_m \ln\!\left(\dfrac{IR}{IR - (v_{th}-v_{rest})}\right)$. Adding the refractory period gives the full inter-spike interval, and its reciprocal is the firing rate:
\begin{eqbox}[LIF F-I curve]
$$f(I) = \left[\, t_{ref} + \tau_m \ln\!\left(\frac{IR}{IR - (v_{th}-v_{rest})}\right)\right]^{-1}, \qquad I > I_{rheo}$$
\end{eqbox}

\begin{keybox}
As $I \to I_{rheo}^+$, the log term diverges $\Rightarrow$ firing rate $\to 0$ continuously (not a jump). As $I \to \infty$, the log term $\to 0$ and $f(I) \to 1/t_{ref}$ — the \textbf{refractory period sets a hard ceiling} on firing rate, and the curve visibly bends over (saturates) as $I$ grows.
\end{keybox}

% =====================================================================
\section{Numerical Integration: Euler vs.\ RK4}
% =====================================================================
Every model here is simulated by discretizing $\dot x = f(x,t)$ in time.

\begin{eqbox}[Forward Euler]
$$x_{n+1} = x_n + \Delta t \cdot f(x_n, t_n)$$
One derivative evaluation per step. Cheap, and the \textbf{project standard for every network simulation}, at $\Delta t \le 0.1$~ms.
\end{eqbox}

\begin{eqbox}[Classical 4th-order Runge-Kutta (RK4)]
$$k_1 = f(x_n),\quad k_2 = f(x_n+\tfrac{\Delta t}{2}k_1),\quad k_3 = f(x_n+\tfrac{\Delta t}{2}k_2),\quad k_4 = f(x_n+\Delta t\,k_3)$$
$$x_{n+1} = x_n + \frac{\Delta t}{6}(k_1 + 2k_2 + 2k_3 + k_4)$$
Four derivative evaluations per step $\Rightarrow$ far more accurate at a given $\Delta t$, but $4\times$ the cost. Used \emph{only} as a single-neuron analytical cross-check, never for networks.
\end{eqbox}

\begin{keybox}[When does the choice actually matter?]
For \textbf{linear} sub-threshold dynamics (Passive, LIF), Euler at $\Delta t=0.1$~ms already tracks RK4 almost exactly — there is little accuracy left on the table. For \textbf{nonlinear} dynamics (EIF, AdEx), Euler's local error grows faster with $\Delta t$: at coarse steps ($\Delta t \gtrsim 2$~ms) Euler starts \emph{missing spikes} that a fine-grained reference or RK4 still detects. This is exactly why the network simulations (1000 neurons $\times$ tens of thousands of steps) stick to cheap Euler at $\Delta t \le 0.1$~ms, where this error is negligible for every model in this project — while RK4 remains available for precision checks on a single neuron.
\end{keybox}

% =====================================================================
\section{Exponential Integrate-and-Fire (EIF)}
% =====================================================================
LIF's hard threshold is a modeling convenience, not a mechanism. Real spike initiation is driven by fast, strongly voltage-dependent Na$^+$ channel activation — better approximated by an exponential acceleration than a step function.

\begin{eqbox}[EIF membrane equation]
$$C_m \frac{\dd v}{\dd t} = -g_L(v - E_L) + g_L \Delta_T \exp\!\left(\frac{v - v_T}{\Delta_T}\right) + I(t)$$
A spike is registered when $v$ reaches a numerical cutoff $v_{peak}$ (standing in for the true, divergent upstroke); then $v \leftarrow v_{reset}$.
\end{eqbox}

\subsection{The role of \texorpdfstring{$\Delta_T$}{Delta\_T} (threshold slope factor)}
\begin{itemize}[itemsep=2pt]
  \item $\Delta_T \to 0$: the exponential term becomes a step at $v_T$ $\Rightarrow$ recovers something close to LIF's hard threshold.
  \item Larger $\Delta_T$: softer, more gradual spike initiation, more visible ``upstroke curvature'' before the peak.
\end{itemize}

\begin{keybox}
Because the exponential term is nonzero even \emph{below} $v_T$, it already contributes a small depolarizing boost pre-threshold. For matched $v_T = v_{th}$, this means \textbf{EIF's effective rheobase is lower than LIF's} — it takes less current to destabilize the resting state once the nonlinearity is doing some of the work.
\end{keybox}

\subsection{Type I excitability}
Both LIF and EIF exhibit \textbf{Type I excitability}: the F-I curve rises continuously from zero at rheobase (no discontinuous jump to a nonzero rate), and the latency to the first spike \textbf{diverges} as $I \to I_{rheo}^+$ — a signature of ``critical slowing down'' near the bifurcation that creates repetitive firing (see \S7).

% =====================================================================
\section{AdEx: Adaptive Exponential Integrate-and-Fire}
% =====================================================================
Neither LIF nor EIF can reproduce something real cortical neurons do constantly: \emph{change their firing rate over the course of a single, constant-current step} (spike-frequency adaptation), or fire in tight bursts. Both require a second state variable — a slow current, $w$, that builds with spiking and feeds back onto $v$.

\begin{eqbox}[AdEx coupled equations]
$$C_m \frac{\dd v}{\dd t} = -g_L(v - E_L) + g_L \Delta_T \exp\!\left(\frac{v - v_T}{\Delta_T}\right) - w + I(t)$$
$$\tau_w \frac{\dd w}{\dd t} = a(v - E_L) - w + b\,\tau_w \sum_f \delta(t - t^f)$$
On each spike ($v \to v_{peak}$): $\quad v \leftarrow v_{reset}, \qquad w \leftarrow w + b$
\end{eqbox}

\begin{defbox}[Reading the $\delta$-function term]
$b\,\tau_w\sum_f \delta(t-t^f)$ means exactly what it looks like once you integrate it: at every spike time $t^f$, $w$ receives an \textbf{instantaneous jump of size $b$}. Between spikes, that term is simply absent, and $w$ relaxes toward $a(v - E_L)$ with time constant $\tau_w$. It is \emph{not} a smoothed or filtered approximation — implementations should apply the jump directly.
\end{defbox}

\subsection{What each parameter controls}
\begin{center}
\begin{tabular}{@{}p{1.6cm}p{2.6cm}p{9.5cm}@{}}
\toprule
\textbf{Param.} & \textbf{Unit} & \textbf{Role} \\
\midrule
$a$ & nS & \emph{Subthreshold} coupling: how strongly $w$ tracks $v$ even without spiking. Negative $a$ (facilitation) is what enables bursting. \\
$b$ & pA & \emph{Spike-triggered} increment to $w$. Larger $b$ = stronger adaptation per spike. \\
$\tau_w$ & ms & Adaptation decay time constant. Small $\tau_w$ = fast recovery between spikes (can allow early bursts); large $\tau_w$ = adaptation lingers, suppressing rate for a long time. \\
$v_{reset}$ & mV & Where $v$ lands after a spike. If close to/above $v_T$, the exponential term is still active post-reset, enabling rapid re-firing (bursts) before $w$ catches up. \\
\bottomrule
\end{tabular}
\end{center}

\subsection{Four regimes, one model}
\begin{center}
\begin{tabular}{@{}lllll@{}}
\toprule
\textbf{Regime} & $a$ (nS) & $b$ (pA) & $\tau_w$ (ms) & $v_{reset}$ (mV) \\
\midrule
Tonic spiking & 0 & 0 & 100 & $-58$ \\
Adapting & 0 & 60 & 200 & $-58$ \\
Initial burst & 0 & 120 & 50 & $-46$ \\
Regular bursting & $-5$ & 100 & 100 & $-46$ \\
\bottomrule
\end{tabular}
\end{center}

\begin{keybox}[Spike-frequency adaptation, quantitatively]
Because AdEx's rate \emph{changes during a trial}, define two F-I curves: the \textbf{instantaneous} rate (from the first ISI, before $w$ has built up) and the \textbf{steady-state} rate (from late ISIs, after $w$ equilibrates). The gap between these two curves, as a function of $I$, \emph{is} spike-frequency adaptation, made visible. LIF has no such gap — its ISIs are constant under constant current.
\end{keybox}

% =====================================================================
\section{Phase-Plane Analysis}
% =====================================================================
The $(v,w)$ phase plane explains \emph{why} the regimes above emerge, geometrically, rather than treating them as arbitrary parameter luck.

\begin{defbox}[Nullclines]
The \textbf{v-nullcline} ($\dot v = 0$) and \textbf{w-nullcline} ($\dot w = 0$):
$$w = -g_L(v-E_L) + g_L\Delta_T e^{(v-v_T)/\Delta_T} + I \qquad\text{(v-nullcline)}$$
$$w = a(v - E_L) \qquad\text{(w-nullcline; a straight line through } (E_L, 0) \text{ with slope } a\text{)}$$
\end{defbox}

\begin{defbox}[Fixed point]
Any $(v,w)$ where \emph{both} nullclines cross: $\dot v = \dot w = 0$ simultaneously. If the trajectory ever sat exactly there, it would stay forever. Away from fixed points, the trajectory flows according to the two derivative equations.
\end{defbox}

\subsection{The mechanism behind spike onset}
At $I=0$, the system typically has \textbf{two} fixed points: a stable one near $E_L$ (rest) and an unstable one further along the $v$-axis (a genuine, geometric threshold — start to its right and the trajectory is pushed toward a spike rather than settling). As $I$ increases, the v-nullcline shifts upward and the two fixed points slide toward each other, until at some critical current they \textbf{collide and annihilate} (a \emph{saddle-node bifurcation}). With no fixed point left to rest at, the trajectory has no choice but to cycle indefinitely: repetitive spiking.

\begin{keybox}
This is \emph{why} a rheobase exists at all, for both EIF and AdEx — it's the current at which the stable fixed point disappears. AdEx's second dimension ($w$) is what lets this same basic mechanism support adaptation and bursting: the $w$-nullcline's slope ($a$) and the reset location relative to the nullclines determine whether the post-spike trajectory lands somewhere that quickly finds a new near-fixed-point (tonic spiking), drifts further away each cycle (adaptation), or periodically re-approaches and escapes a transient fixed point (bursting).
\end{keybox}

% =====================================================================
\section{Synaptic Dynamics}
% =====================================================================
\subsection{Exponential synapse}
\begin{eqbox}[Exponential kernel]
$$\tau_{syn}\frac{\dd g}{\dd t} = -g, \qquad g \leftarrow g + w_{syn} \text{ on each incoming spike}$$
Closed-form response to a single impulse at $t=0$: $g(t) = w_{syn}\,e^{-t/\tau_{syn}}$.
\end{eqbox}

\subsection{Alpha synapse}
A smoother rise-and-decay kernel, built from two coupled linear ODEs (both with the same $\tau$):
$$\frac{\dd x}{\dd t} = -\frac{x}{\tau} \quad (x \leftarrow x + w_{syn}\text{ per spike}), \qquad \frac{\dd g}{\dd t} = \frac{x - g}{\tau}$$
The resulting $g(t)$ peaks at exactly $t=\tau$ after an isolated input spike (Rall's alpha function).

\subsection{Poisson background input}
\begin{defbox}[Bernoulli-per-bin approximation]
At each timestep $\Delta t$ (ms), a source with rate $R$ (Hz) spikes independently with probability
$$p_{spike} = \frac{R\cdot\Delta t}{1000}$$
This is an excellent approximation to a true Poisson process whenever $R\Delta t \ll 1000$ — comfortably true at $\Delta t \le 0.1$~ms for any biologically plausible rate.
\end{defbox}

\begin{keybox}[Shot-noise statistics — why this matters for network tuning]
Filtering a Poisson spike train (rate $R$, in spikes/ms) through an exponential synapse (weight $w$, decay $\tau$) produces a current with
$$\text{mean} = R\,w\,\tau, \qquad \text{variance} \approx \frac{R\,w^2\,\tau}{2}$$
For a \emph{fixed mean}, increasing $w$ while decreasing $R$ proportionally \textbf{increases the variance} (fewer, larger events = more fluctuation). This is the knob that separates \textbf{mean-driven} firing (large mean current, pushes steadily over threshold, fairly regular ISIs) from \textbf{fluctuation-driven} firing (mean held below rheobase, but variance large enough that fluctuations occasionally cross threshold — irregular ISIs, CV closer to 1). The balanced network relies on being in the fluctuation-driven regime.
\end{keybox}

% =====================================================================
\section{Balanced Recurrent Networks}
% =====================================================================
Scaling from one neuron to $N$ (here, 1000) sparsely-connected LIF neurons.

\begin{defbox}[Dale's law]
Each neuron's \emph{outgoing} synaptic weights are either all excitatory (positive) or all inhibitory (negative) — never mixed, matching the biological constraint that a real neuron releases one neurotransmitter type. In this project, the first 80\% of neuron indices are excitatory, the remaining 20\% inhibitory.
\end{defbox}

\subsection{Network construction}
\begin{itemize}[itemsep=2pt]
  \item Connection probability $p$: any ordered pair $(i,j)$ is connected with probability $p$, independently $\Rightarrow$ mean in-degree $\approx N\!\cdot\!p$, approximately binomially distributed.
  \item Inhibitory weights scaled: $w_{inh} = -g \cdot w_{exc}$, with $g$ typically $\gg 1$ (e.g., $g=5$) — inhibition needs to be \emph{stronger, not just present}, per connection.
  \item Every neuron also receives independent Poisson background drive (\S8.3) — without it, a sparsely-connected, purely leaky network would simply stay silent forever.
\end{itemize}

\subsection{The Asynchronous-Irregular (AI) regime}
\begin{defbox}[AI regime — three signatures]
\begin{enumerate}[itemsep=2pt, leftmargin=1.4em]
  \item \textbf{Asynchronous}: no shared population-wide rhythm; raster plot shows no vertical stripes/bursts; population rate trace is flat, not oscillatory.
  \item \textbf{Irregular}: individual spike trains resemble a memoryless (Poisson) process. Measured via the \textbf{ISI coefficient of variation}:
  $$CV = \frac{\text{std}(\text{ISIs})}{\text{mean}(\text{ISIs})}, \qquad CV\approx 0 \text{ (clock-like)}, \; CV\approx 1 \text{ (Poisson-like)}, \; CV>1 \text{ (bursty)}$$
  \item \textbf{Stable, low rate}: mean population rate settles to a low value and stays there, rather than diverging (runaway excitation) or vanishing (net inhibition dominates completely).
\end{enumerate}
\end{defbox}

\subsection{Population rate}
Binning all spikes network-wide into windows of width $\Delta_{bin}$ (seconds):
$$r(t) = \frac{\text{spike count in bin}}{N \cdot \Delta_{bin}} \quad \text{[Hz]}$$

\subsection{Why ``balanced''? The mechanism, not just the name}
\begin{keybox}
With \textbf{weak} inhibition (small $g$), increasing excitatory weight directly raises the population rate — recurrent excitation amplifies itself largely unchecked (approaching runaway/mean-driven behavior). With \textbf{strong} inhibition (large $g$), the population rate becomes largely \emph{insensitive} to further increases in excitatory weight: inhibition tracks and cancels the extra excitatory drive almost as fast as it arrives. That insensitivity — not any particular numerical rate — is the actual definition of ``balance,'' and it's what keeps the net input to each neuron fluctuation-dominated (small, noisy net current) rather than mean-dominated (large, steady net current), which is in turn what produces the irregular, asynchronous firing described above.
\end{keybox}

% =====================================================================
\newpage
\section{Quick-Reference Formula Sheet}
% =====================================================================
\begin{center}
\renewcommand{\arraystretch}{1.35}
\begin{longtable}{@{}p{3.6cm}p{12.5cm}@{}}
\toprule
\textbf{Concept} & \textbf{Formula} \\
\midrule
\endhead
Passive membrane & $\tau_m \dot v = -(v-v_{rest}) + RI(t)$, \quad $\tau_m = C_m/g_L$ \\
Passive step response & $v(t) = v_{rest} + \dfrac{I}{g_L}\left(1-e^{-t/\tau_m}\right)$ \\
LIF rule & $v\!\ge\! v_{th} \Rightarrow v\!\leftarrow\! v_{reset}$, refractory $t_{ref}$ \\
LIF rheobase & $I_{rheo} = g_L(v_{th}-v_{rest})$ \\
LIF F-I curve & $f(I) = \left[t_{ref} + \tau_m \ln\!\frac{IR}{IR-(v_{th}-v_{rest})}\right]^{-1}$ \\
EIF & $C_m \dot v = -g_L(v-E_L) + g_L\Delta_T e^{(v-v_T)/\Delta_T} + I$ \\
AdEx ($v$) & $C_m \dot v = -g_L(v-E_L) + g_L\Delta_T e^{(v-v_T)/\Delta_T} - w + I$ \\
AdEx ($w$) & $\tau_w \dot w = a(v-E_L) - w + b\tau_w\sum_f \delta(t-t^f)$ \\
AdEx spike reset & $v\!\leftarrow\! v_{reset}$, \quad $w \leftarrow w + b$ \\
v-nullcline & $w = -g_L(v-E_L) + g_L\Delta_T e^{(v-v_T)/\Delta_T} + I$ \\
w-nullcline & $w = a(v-E_L)$ \\
Euler step & $x_{n+1} = x_n + \Delta t\, f(x_n)$ \\
RK4 step & $x_{n+1} = x_n + \frac{\Delta t}{6}(k_1+2k_2+2k_3+k_4)$ \\
Exponential synapse & $\tau_{syn}\dot g = -g$, \; $g \leftarrow g+w_{syn}$ per spike \\
Alpha synapse & $\dot x = -x/\tau$ (jump $w_{syn}$ per spike), \; $\dot g = (x-g)/\tau$ \\
Poisson bin probability & $p_{spike} = R\Delta t / 1000$ \\
Shot-noise mean/variance & mean $=Rw\tau$, \; variance $\approx Rw^2\tau/2$ \\
ISI coefficient of variation & $CV = \text{std(ISI)}/\text{mean(ISI)}$ \\
Population rate & $r(t) = \dfrac{\text{spikes in bin}}{N\cdot\Delta_{bin}\text{[s]}}$ \\
Dale's law weight scaling & $w_{inh} = -g\cdot w_{exc}$, \; $g \gg 1$ for balance \\
\bottomrule
\end{longtable}
\end{center}

	\newpage
% =====================================================================
\section{Self-Test Questions}
% =====================================================================
\begin{enumerate}[itemsep=4pt, leftmargin=1.6em]
  \item Why can a passive membrane never produce a spike, regardless of input current magnitude?
  \item Derive (or explain in words) why the LIF F-I curve's firing rate approaches $1/t_{ref}$ as $I \to \infty$.
  \item Why does EIF have a \emph{lower} rheobase than an LIF neuron with the same $v_T = v_{th}$?
  \item What does $\Delta_T \to 0$ do to the EIF model, intuitively?
  \item Write out what the term $b\,\tau_w\sum_f \delta(t-t^f)$ actually does to $w$ at a spike time — in one sentence, without integral notation.
  \item Why does AdEx need \emph{two} state variables to produce bursting, when LIF/EIF (one variable) cannot?
  \item In the $(v,w)$ phase plane, what does it mean, geometrically, for a neuron to be at rest? What happens as input current pushes the system toward its rheobase?
  \item What is the difference between the "instantaneous" and "steady-state" F-I curve, and why does only AdEx (not LIF) have a gap between them?
  \item Why does a Poisson background source need a nonzero rate \emph{and} a large enough synaptic weight to actually drive spiking, rather than just a high rate alone?
  \item Explain, in terms of mean vs.\ variance of shot noise, the difference between a "mean-driven" and a "fluctuation-driven" spiking regime.
  \item What does Dale's law constrain, specifically, in the connectivity matrix $W$?
  \item Why is the Asynchronous-Irregular regime described as an \emph{emergent network property} rather than a property of any single neuron?
  \item If you doubled every excitatory weight in a balanced network \emph{without} increasing inhibition to match, what would you expect to happen to the population rate? What if inhibition were already very strong before you doubled excitation?
  \item Why does this project use Euler for every network simulation but keep RK4 available only for single neurons?
\end{enumerate}

	\newpage
\vspace{0.3cm}
\begin{center}\rule{0.4\textwidth}{0.4pt}\end{center}
\subsection*{Answers (brief)}
{\small
\begin{enumerate}[itemsep=2pt, leftmargin=1.6em]
  \item There is no threshold/reset rule in the model at all — $v$ only asymptotes toward $v_{rest}+IR$; nothing triggers a spike event no matter how high that asymptote is.
  \item As $I\to\infty$, the charging time from reset to threshold $\to 0$ (the log term vanishes), leaving only the refractory period between spikes, so $f \to 1/t_{ref}$.
  \item The exponential current is already nonzero (and depolarizing) below $v_T$, effectively pre-loading some of the work a hard threshold would otherwise require more current to do.
  \item It sharpens the exponential term into something close to a step at $v_T$, making EIF behave increasingly like LIF.
  \item At each spike, $w$ jumps up instantly by exactly $b$; between spikes, that term contributes nothing and $w$ just relaxes toward $a(v-E_L)$.
  \item A single variable resets to the same state after every spike, so every inter-spike interval under constant input would look identical; $w$ carries memory of recent spiking across intervals, letting successive ISIs differ (lengthen, shorten, or cluster).
  \item Rest = sitting at a stable fixed point (both nullclines cross, flow points inward). As $I$ rises, the v-nullcline shifts and the stable/unstable fixed-point pair merges and vanishes (saddle-node bifurcation) — with nowhere to rest, the system is forced into a repetitive spiking limit cycle.
  \item Instantaneous = rate from the very first ISI (before $w$ has built up); steady-state = rate from late ISIs (after $w$ has equilibrated). LIF has no adaptation variable, so its rate never changes within a trial — no gap.
  \item Rate alone sets how \emph{often} impulses arrive; weight sets how \emph{large} each one is. A high rate of vanishingly small impulses can still average out to zero net depolarization if the weight is too small to matter individually or in sum relative to threshold.
  \item Mean-driven: the average current alone is enough to push $v$ steadily up to threshold nearly every cycle $\Rightarrow$ fairly regular ISIs. Fluctuation-driven: the mean sits below threshold, and only the random fluctuations (variance) occasionally push $v$ over $\Rightarrow$ irregular, Poisson-like ISIs.
  \item Every \emph{column} of $W$ (one presynaptic neuron's outgoing weights) must be entirely $\ge 0$ (excitatory source) or entirely $\le 0$ (inhibitory source) — never a mix of signs within one column.
  \item No single LIF neuron in the network does anything qualitatively different from an isolated LIF neuron under noisy current injection — the asynchronous, irregular, stable-rate \emph{population} behavior only exists because of how 1000 of them are wired together (Dale's law + sparse connectivity + inhibition scaled above excitation).
  \item Weak inhibition to begin with: rate rises noticeably (amplification). Already-strong inhibition: rate stays roughly the same — inhibition absorbs the extra drive, which is exactly the balance mechanism.
  \item Networks simulate many neurons over many timesteps, where Euler's cost (1 derivative evaluation/step) is what makes the simulation tractable at all; RK4's $4\times$ cost is only worth paying when precisely validating one neuron's trajectory against a reference.
\end{enumerate}
}
	\newpage
% =====================================================================
\section{Glossary}
% =====================================================================
\begin{description}[style=nextline, leftmargin=1.2em, itemsep=3pt]
  \item[Rheobase] Minimum constant current for sustained/eventual firing.
  \item[Refractory period ($t_{ref}$)] Time after a spike during which the neuron cannot fire again; sets the maximum possible firing rate to $1/t_{ref}$.
  \item[Nullcline] The curve in a phase plane along which one state variable's derivative is zero.
  \item[Fixed point] A point where \emph{every} state variable's derivative is zero simultaneously.
  \item[Bifurcation] A qualitative change in a dynamical system's behavior (e.g., number of fixed points) as a parameter crosses a critical value.
  \item[Saddle-node bifurcation] Two fixed points (one stable, one unstable) collide and annihilate as a parameter changes — the mechanism behind spike onset in EIF/AdEx.
  \item[Type I excitability] Firing rate rises continuously from zero at rheobase (as opposed to Type II, which jumps to a nonzero rate).
  \item[Spike-frequency adaptation (SFA)] Progressive lengthening of ISIs over a spike train under constant input.
  \item[Dale's law] A neuron's outgoing synapses are either all-excitatory or all-inhibitory, never mixed.
  \item[Coefficient of variation (CV)] std/mean of inter-spike intervals; a standard irregularity measure.
  \item[Asynchronous-Irregular (AI) regime] Network state with no shared rhythm, individually irregular spiking, and a stable low mean rate.
  \item[Mean-driven vs.\ fluctuation-driven] Whether spiking is caused mainly by the average input current crossing threshold, or by noise/variance occasionally pushing a subthreshold-on-average current over threshold.
  \item[Shot noise] The current fluctuations produced by filtering a discrete (Poisson) spike train through a continuous synaptic kernel.
\end{description}

\end{document}
